\chapter{Introducción para los jugadores}

\begin{CajaSecreto}
\lead{Nota para el Director:} lee el texto de apertura a tu mesa. Establece un tono asfixiante, urbano y noir. Los personajes no están aún en los abismos de la fundición de 2102; están en «La Margen» de 2084, rodeados de asfalto mojado, neones sucios, burocracia opresiva y lluvia ácida. Bajo ningún concepto reveles el verdadero propósito del Proyecto IZAR antes de la Fase 1.
\end{CajaSecreto}

\section{Lectura para la mesa: El espejismo del mérito}

\neofigure{01_distrito7_la_margen.jpg}{Distrito 7: el sirimiri ácido disuelve el hollín sobre el asfalto de La Margen.}

\begin{quote}
«El sirimiri en el Distrito 7 nunca limpia las calles; solo esparce el óxido y difumina los neones de las pantallas publicitarias sobre el asfalto resbaladizo. Bienvenidos a «La Margen», el cinturón urbano de clase media-baja de NeoBilbao. Aquí no sois esclavos de la fundición profunda; sois algo peor: sois ciudadanos con esperanza.

Vivís en micro-apartamentos cápsula de catorce metros cuadrados, pagáis vuestros créditos de oxígeno a plazos mensuales y trabajáis como técnicos de mantenimiento, celadores de hospital u oficinistas de archivo en un sistema que os exprime la vida hora tras hora. Pero aguantáis. Todos en este barrio aguantan por culpa del Programa Azul de Excelencia Ciudadana.

En las tascas oscuras y bajo los paraguas de polímero transparente del tren mag-lev, todos hablan de lo mismo: si eres un ciudadano ejemplar, si sacas buenas notas en tus evaluaciones y mantienes un Índice Azul impecable, la corporación te selecciona. Te otorgan un piso en la Cúpula Dorada, empleo ejecutivo y aire puro para siempre. \emph{«El talento tiene su recompensa»}, reza el eslogan que parpadea en las pantallas CRT de Orekak-Corp.

Hace seis meses, Ane Garate, una brillante ingeniera de telecomunicaciones de vuestro propio bloque, ganó el Programa Azul. El barrio entero salió a celebrarlo con kalimotxo sintético. Fue el orgullo de la vecindad.

Esta noche, su familia os ha citado de urgencia en una tasca clandestina bajo las vías del mag-lev. Ane ha vuelto. Físicamente parece ilesa, pero no deja de temblar, habla de sí misma en plural y tiene la mirada congelada de quien ha contemplado el vacío absoluto.

Pedís un trago mientras el sirimiri golpea con furia el cristal blindado. Algo huele a podrido en la lotería de NeoBilbao, y vosotros estáis a punto de descubrir el porqué.»
\end{quote}

\section{El reencuentro en la Tasca de Jon Ander}

\neofigure{02_ane_garate_reaparicion.jpg}{Ane Garate en el reservado de la tasca: pupilas dilatadas y un tosco puerto en la nuca.}

La escena arranca en la trastienda del bar «Zulo», regentado por Jon Ander. Ane está sentada en un rincón con un chaquetón empapado. 

\begin{itemize}
  \item \lead{El síntoma físico:} Bajo su pelo rapado en la coronilla, asoma un puerto de conexión tosco, de titanio de grado industrial, asegurado con cuatro tornillos quirúrgicos que supuran un líquido amarillento. No es wetware comercial estético; parece el anclaje de un transformador de alta tensión.
  \item \lead{El habla fragmentada:} Ane murmura: \emph{«Éramos dieciséis en el nodo... El bus de datos ardía. No desconectéis el bus... el perro negro nos huele...»}. Al miraros, reconoce vuestros rostros solo por destellos.
  \item \lead{La baliza biofísica (Mejora de tensión):} El puerto de Ane emite un zumbido electromagnético que hace oscilar los tubos fluorescentes de la tasca y satura las radios de onda corta. Si los personajes no aíslan la señal o no la trasladan con rapidez, los Auditores corporativos no tardarán en triangular su posición.
\end{itemize}

\section{Prólogo audiovisual: El eco de la Lotería}

Guion cinemático de un minuto y quince segundos, diseñado para ser narrado o reproducido antes de iniciar la primera sesión.

\begin{CajaCinematica}{Ficha del tráiler}
\cine{Título:} NEOBILBAO 2084: La Lotería Azul\par
\cine{Duración aproximada:} 1 minuto y 15 segundos\par
\cine{Tono visual:} ciberpunk noir, lluvia ácida constante, pantallas CRT parpadeantes en fósforo verde y ámbar, gabardinas mojadas, reflejos de neón cian en charcos de aceite.\par
\cine{Música y sonido:} sintetizadores analógicos pesados estilo Vangelis/Carpenter, zumbido de transformadores eléctricos, el traqueteo sordo del tren mag-lev y el goteo corrosivo del sirimiri.
\end{CajaCinematica}

\begin{CajaCinematica}{Escena 1 -- La ilusión del mérito (0:00 -- 0:25)}
\cine{Visual:} plano general de La Margen al anochecer. Kilómetros de bloques colmena de hormigón gris cruzados por cables coaxiales colgantes. En lo alto de una torre, una gigantesca pantalla publicitaria CRT muestra a una familia sonriente de traje blanco entrando en la Cúpula Dorada bajo el lema: \textbf{OREKAK-CORP: EL TALENTO TIENE SU RECOMPENSA}.\par
\cine{Voz en off:}\par
\textit{«Nos dijeron que el mérito era la única llave. Que si estudiabas catorce horas diarias, si no protestabas y mantenías limpio tu expediente, la Lotería Azul te sacaría del barro. Todos nos lo creímos. Todos queríamos creer».}
\end{CajaCinematica}

\begin{CajaCinematica}{Escena 2 -- El regreso del fantasma (0:25 -- 0:50)}
\cine{Visual:} la cámara desciende a ras de suelo entre paraguas transparentes y humo de alcantarilla. Corte rápido a la trastienda de una tasca. La luz mortecina de una lámpara ilumina la nuca de Ane Garate: una interfaz hexagonal tosca atornillada al cráneo, rodeada de costras y cables de cobre expuestos. Sus manos tiemblan sobre un vaso de plástico.\par
\cine{Efectos de sonido:} un chasquido agudo de estática electromagnética que interrumpe el audio ambiental.\par
\cine{Voz en off (con eco desfasado y varias voces simultáneas):}\par
\textit{«No es un premio... Nunca fue un premio. No hay pisos en la Cúpula. Solo hay tanques... cables gruesos... y frío. Mucho frío».}
\end{CajaCinematica}

\begin{CajaCinematica}{Escena 3 -- Los Secretas en la niebla (0:50 -- 1:15)}
\cine{Visual:} en el exterior de la tasca, una furgoneta blanca frena sin hacer ruido en el callejón. Dos figuras con gabardinas largas y gafas de sol oscuras estilo 90s bajan del vehículo desenfundando inyectores neumáticos de toxina. La cámara enfoca la pantalla de un terminal portátil donde el pulso biométrico de Ane parpadea en rojo.\par
\cine{Texto en pantalla:}\par
\textbf{NEOBILBAO 2084 // LA LOTERÍA AZUL}\par
\textbf{El silencio tiene un precio.}\par
\cine{Efectos de sonido:} el siseo presurizado de un inyector neumático cargándose y un trueno sordo sobre la Ría.
\end{CajaCinematica}
